\documentclass[twoside]{article}
\usepackage[a5paper,margin=0.5in]{geometry}
\usepackage{gregosheet}
\usepackage{gregosheet-psalm}
\usepackage{lyluatex}
% Font sizes
\newcommand{\musicfontsize}{20}
\newcommand{\lyricfontsize}{10}
\newcommand{\psalmfontsize}{10}

% Fonts
% Change this to any font command you want:
% \rmfamily (roman/serif), \sffamily (sans-serif), \ttfamily (monospace)
% or use \setmainfont{FontName} or \newfontfamily for specific fonts
\newfontfamily\lyricfont{Cambria}
\newfontfamily\psalmfont[BoldFont={Cambria Bold}]{Cambria}
\newfontfamily\headingfont[BoldFont={Cambria Bold}]{Cambria}
\setmainfont{Cambria}

% Spacing
\newcommand{\systemvskip}{10pt}
\newcommand{\lyricvskip}{2pt}
\newcommand{\blockvskip}{10pt}

% Page layout
\usepackage{fancyhdr}
\pagestyle{fancy}
\fancyhf{}
\fancyhead[CO]{\textcolor{red}{\MakeUppercase{\leftmark}}}
\fancyhead[RO]{\thepage}
\fancyhead[LE]{\thepage}
\fancyhead[CE]{\textcolor{red}{\MakeUppercase{\rightmark}}}
\renewcommand{\headrulewidth}{0pt}
\setlength{\parindent}{0pt}

% Title style
\makeatletter
\renewcommand{\maketitle}{
  \begin{titlepage}
    \centering
    \vspace*{\fill}
    {\headingfont\fontsize{32}{38}\selectfont\bfseries\MakeUppercase{\@title}\par}
    \vskip 2em
    {\headingfont\fontsize{16}{20}\selectfont\MakeUppercase{\@subtitle}\par}
    \vspace*{\fill}
  \end{titlepage}
}
\newcommand{\subtitle}[1]{\gdef\@subtitle{#1}}
\newcommand{\@subtitle}{}
\makeatother

% Section style
\usepackage{titlesec}
\titleformat{\section}
  {\centering\headingfont\fontsize{20}{24}\selectfont\bfseries}
  {}{0pt}{\MakeUppercase}
\titlespacing*{\section}{0pt}{\blockvskip}{\blockvskip}

% Subsection style
\titleformat{\subsection}
  {\centering\headingfont\fontsize{16}{20}\selectfont\bfseries\addfontfeature{LetterSpace=10.0}}
  {}{0pt}{\MakeUppercase}
\titlespacing*{\subsection}{0pt}{\blockvskip}{\blockvskip}

% Subsubsection style
\titleformat{\subsubsection}
  {\centering\headingfont\fontsize{12}{14}\selectfont\bfseries\scshape}
  {}{0pt}{}
\titlespacing*{\subsubsection}{0pt}{\blockvskip}{\blockvskip}

% Part title command
\newcommand{\parttitle}[2][]{
  \vskip\blockvskip
  \par\noindent
  \headingfont\fontsize{10}{12}\selectfont
  \textcolor{red}{\MakeUppercase{#2}}%
  \ifx&#1&\else\hfill\textcolor{red}{#1}\fi
  \nopagebreak
  \par
}

% Rubric command
\newcommand{\rubric}[1]{
  \par\noindent
  {\fontsize{8}{10}\selectfont\textcolor{red}{#1}}
  \par
}

% Red text command
\newcommand{\redtext}[1]{
  \par\noindent
  \textcolor{red}{#1}
  \par
}

\newcommand{\p}{\textcolor{red}{P.: }}

\newcommand{\h}{\textcolor{red}{H.: }}

\newcommand{\e}{\textcolor{red}{E.: }}

\newcommand{\lec}{\textcolor{red}{L.: }}

\newcommand{\diac}{\textcolor{red}{D.: }}

% Red dash command
\newcommand{\reddash}{\textcolor{red}{\textemdash{}}\xspace}
\usepackage{xspace}

% Lectio environment
\newenvironment{lectio}[1]{
  \par\noindent
  {\centering\textcolor{red}{\textit{#1}}\par}
  \setlength{\parindent}{0.2in}
  \everypar{\setlength{\parindent}{0.2in}}
}{
  \par
}

% Versicle-Response command
\newcommand{\versicle}[2]{
  \par\noindent
  \hangindent=1em\hangafter=1
  \emergencystretch=1em
  \makebox[1em][l]{\textcolor{red}{V.}}#1
  \par\noindent
  \hangindent=1em\hangafter=1
  \emergencystretch=1em
  \makebox[1em][l]{\textcolor{red}{F.}}#2
  \par
}

% Hymn stanzas environment
\usepackage{multicol}
\newenvironment{hymnstanzas}{
  \setlength{\leftmargin}{0pt}
  \setlength{\rightmargin}{0pt}
  \setlength{\columnsep}{4em}
  \begin{multicols}{2}
  \raggedright
}{
  \end{multicols}
}

\newcommand{\stanzabreak}{%
  \par\addvspace{0.8\baselineskip}%
}


% Single centered stanza (for odd last stanza)
\newcommand{\centerstanza}[2][0.4]{
  \par
  \begin{center}
  \begin{minipage}{#1\textwidth}
  \raggedright
  #2
  \end{minipage}
  \end{center}
  \par
}

% Structured text environment
\newenvironment{preces}[2]{
  \par\noindent
  #1
  \par\noindent
  \hskip0.2in\textit{#2}
  \par
  \setlength{\parindent}{0pt}
  \everypar{\hangindent=0.2in\hangafter=1}
}{
  \par
}

\newcommand{\printsheet}[1]{%
  \ifcsname sheet@#1\endcsname
    \csname sheet@#1\endcsname
  \else
    \textbf{[Missing sheet: #1]}%
  \fi
}

\usepackage{ragged2e}


\title{Nagyhét 2026}
\subtitle{A szentendrei ferences kápolna szertartásrendje}

\begin{document}

\maketitle

\section{Virágvasárnap}

\commentary{Ezen a napon ünnepli az Egyház Krisztus Urunk Jeruzsálembe való bevonulását a húsvéti %
misztérium beteljesítésére.}

\subsection{Körmenet}

\subsubsection{I. stáció: Az ágak megáldása}

\commentary{A pap és a segédkezők már a szentmiséhez felöltözve a kápolna végébe %
vonulnak. Eközben a következő antifónát énekeljük:}

\parttitle[Hosanna Filio David]{Anttifóna}

\begin{gregosheet}[tone=7. lá tónus]
  \melody{<ô-1--1t---5--:-5--3---4--5z-5--:-5zZ5-4-:-4-4--4--4--4--2--3---1-,-0--2--4--4---4--4t-rR3-:-1--1t---5x-4-2--3---1---1--,,-3-4-[-4-5-6-†-4-,-[-6-5-4-†-rR3-2-.}
  \lyrics{   Ho-zsan-na * Dá-vid Fi-á--nak, ál-dott, a-ki az Úr ne-vé-ben jő, Iz-ra-el-nek Ki-rá-lya,  Ho-zsan-na a ma-gas-ság-ban.}
\end{gregosheet}

\commentary{Ha az idő engedi, a 117. zsoltár verseivel énekeljük:}

\begin{psalmsheet}[tone=7, continuous]
  \psalmtext{
    1. Hálát adjatok az Úrnak, mert jó, * mert az ő irgalmassága örökké való. \\
    2. A Kő, amelyet megvetettek az építők, * az lett most Szegletkővé. ANT. \\
    3. Az Úrnak műve ez * és csodálatos a mi szemeink előtt. \\
    4. Tartsatok ünnepeket lombos ágakkal * föl, egészen a magas oltárig. ANT.
}
\end{psalmsheet}

\commentary{A pap a szokott módon üdvözli a népet (ahogyan a mise elején szokta); rövid intelmet %
ad, hogy a hívek tevékenyen és tudatosan vegyenek részt ennek a napnak a megünneplésében. Ezekkel %
vagy hasonló szavakkal fordul hozzájuk:}

\p Kedves Testvéreim! A Nagyböjt eleje óta szívünket bűnbánattal és jócselekedetekkel már
előkészítettük. Ma azért gyűltünk össze, hogy az egész Egyházzal együtt lélekben előre átéljük
Urunk húsvéti misztériumát vagyis szenvedését és feltámadá\-sát. Ő ugyanis ezt a misztériumot akarta
beteljesíteni, amikor fölment városába, Jeruzsálembe. Szívünk élő hitével és odaadásával emlékezünk
meg erről az üdvösségszerző bevonulásról. Lépjünk mi is az Úr nyomába, hogy amint kegyelméből
társai vagyunk a kereszthordozásban, ugyanúgy feltámadásának és örök életének is részesei lehessünk.

\parttitle{Könyörgés}

\commentary{Az intelem után a pap a következő két könyörgés egyikét imádkozza összetett kézzel:}

\p Könyörögjünk! Mindenható, örök Isten, szenteld meg áldásoddal \textcolor{red}{†} ezeket az ágakat,
hogy mi, akik most örvendezve lépünk Krisztus Király nyomába, őáltala egykor az örök Jeruzsálembe
is eljussunk. Aki él és uralkodik mindörökkön-örökké.

\h Ámen.

\commentary{Vagy:}

\p Könyörögjünk! Urunk, Istenünk, növeld azoknak hitét, akik benned remélnek. Hallgasd meg könyörgő
imánkat, hogy mi, akik ma virágzó ágakat hozunk a diadalmasan bevonuló Krisztus Király elé, a
jócselekedetek gyümölcsét teremjük őáltala. Aki él és uralkodik mindörökkön-örökké.

\h Ámen.

\commentary{A pap szenteltvízzel meghinti az ágakat, minden kísérő szöveg nélkül.}

\parttitle[Mt 21,1-11]{Evangélium}
\p Az Úr legyen veletek!

\h És a te l\underline{e}lkeddel!

\p Evangélium Szent Mát\underline{é} könyvéből.

\h Dicsőség nek\underline{e}d, Istenünk!

\p Amikor Jeruzsálemhez közeledve az Olajfák-hegyére, Betfagéba értek, Jézus el\-küldte két tanítványát
ezekkel a szavakkal: „Menjetek előre a szemközti faluba. Ott mindjárt találni fogtok egy szamarat
megkötve és vele csikóját. Oldjátok el és vezessétek hozzám! Ha valaki szólna valamit, mondjátok,
hogy az Úrnak van rá szüksége, és mindjárt elengedi őket.”

Ez azért történt, hogy beteljesedjék, amit a próféta jövendölt: „Mondjátok meg Sion
lányának: Íme, a királyod érkezik hozzád, szelíden, szamárháton ülve, egy teherhor\-dó
állat csikóján.”

A tanítványok elmentek s úgy tettek, ahogy Jézus meghagyta nekik. Elhozták a szamarat és a
csikóját, letakarták ruháikkal, ő pedig felült rá.

A tömegből nagyon sokan az útra terítették ruháikat, mások ágakat tördeltek a fák\-ról
és az útra szórták. Az előtte járó és az utána vonuló tömeg így kiáltott: „Hozsanna Dávid fiának!
Áldott, aki az Úr nevében jön! Hozsanna a magasságban!”

Amikor beért Jeruzsálembe, megmozdult az egész város, és kérdezgették: „Ki ez?” A
tömeg ezt felelte: „Ő a Próféta, Jézus, a galileai Názáretből.”

\p Ezek az evangéli\underline{u}m igéi.

\h Áldunk tég\underline{e}d, Krisztus!

\parttitle{Homília}

\commentary{Az evangélium után, ha alkalmas, a pap rövid homíliát mondhat. A körmenet elindulására
vagy maga a pap, vagy más erre alkalmas segédkező adja meg a jelt, ezekkel vagy hasonló szavakkal:}
\p Kedves Testvéreim! Kövessük a Jézust ünneplő jeruzsálemi nép példáját, és induljunk békével!

\h Krisztus nevében, ámen.

\commentary{A körmenet elindul a szentély felé.}

\subsubsection{II. stáció: Az ágak leterítése}

\commentary{A körmenet egy alkalmas helyen megáll, a pap átveszi a körmeneti keresztet, miközben
gyermekek lépnek elébe, ágakat emelnek a magasba, majd leterítik azokat a körmenet elé a földre.
A pap a kereszttel kezében áthalad az ágak fölött. Eközben az előénekesek a következő antifónát
intonáják.}

\parttitle[Pueri hebraeorum portantes]{Antifóna}

\begin{gregosheet}
  \melody{<X-1--3--1---0---3----4---3t-5--:-5---5--5u--5-!-5-4--3--4t--5-:-5-4--5--3----4--eE2-1--,-3----eE2-1---0--0---1---3r--3--,-3--4t---tT4-!-2-3--2---1---1-.}
  \lyrics{   Je-ru-zsá-lem gyer-mek né-pe * pál-ma-fák-nak á-ga-it hoz-ván e-lé-be ment az Úr--nak, s_ki-ál--tot-ta han-gos szó-val: Ho-zsan-na    a ma-gas-ság-ban!}
\end{gregosheet}

\subsubsection{III. stáció: A ruhák leterítése}

\commentary{Gyermekek díszes ruhát hoznak és leterítik az útra. A pap, ismét kereszttel kezében,
áthalad rajtuk. Eközben az előénekesek a következő antifónát intonálják:}

\parttitle[Pueri hebraeorum vestimenta]{Antifóna}

\begin{gregosheet}
  \melody{<X-1--3--1---0---3----4---3t-5--:-5--5---5u-5-!-5--5--4--3--4t--5-:-5----4---5---3---4--eE2-1---,-3--eE2--qQ0-§-0--1---3--3r--3-,-3--4-----5--tT4-!-2--3--2--1--1-.}
  \lyrics{   Je-ru-zsá-lem gyer-mek né-pe * le-ter-ít-vén ru-há-it az út-ra,  nagy szó-val így ki-ál--tott: Ho-zsan-na    Dá-vid fi-á-nak,  ál-dott, ki jő    az Úr ne-vé-ben!}
\end{gregosheet}

\subsubsection{IV. stáció: Hódolat Krisztus Király előtt}

\commentary{A körmenet az szentély lépcsőjénél megáll, a pap a keresztet kezébe veszi, s a kereszt
elé járulnak a verzusok énekesei (lehetőség szerint gyerekek). Az éneket ők kezdik, s mindnyájan
megismételjük a refrént. Ezután éneklik a verseket. A versek után közösen megismételjük a refrént.}

\parttitle[Gloria, laus]{Himnusz}

\begin{gregosheet}
  \melody{<X-5--rR3-4y--3--4--4---5---3--2--1--:-3---3---1--3--3--2y---,-3t-5x--5----4--3--4----5--5x--3---2----1-0---eE2-1--1-.}
  \lyrics{   Di-cső-ség és di-csé-ret Te-né-ked, Meg-vál-tó Ki-rá-lyunk, ki-nek gyer-me-ki szép se-reg mon-dott é-kes é---ne-ket.}
\end{gregosheet}

\commentary{Verzusok:}

\begin{gregosheet}
  \melody{<X-1--[----------5u-4--,-7--7---8---7---™6--5-4----5---5-,-5--5--5--5--5u-4-,-7----7---8-----7----™6-5---4--5--5-----,,---1--5----5--5--5---5-5u----4-,-7----7---8---7--™6---5---4---5--5--,-5----5---[---------5u-4--,-7--7---8---7--™6---5--4----5--5-----,,---1-5---5---5-5u-4-,-7-7-8-7-™6-5-4-5-5-,-[-5u-4-,-7-7-8-7-™6-5-4-5-5-,,-1-[-5u-4-,-7-7-8-7-™6-5-4-5-5-,-[-5u-4-,-7-7-8-7-™6-5-4-5-5-,,-1-5-5-5-5-5u-4-,-7-7-8-7-™6-5-4-5-5-,-[-5u-4-,-7-7-8-7-™6-5-4-5-5-.}
  \lyrics{   Iz-raelnek_Ki-rá-lya: Dá-vid-nak rég-től i-gért sar-ja, ki az Úr ne-vé-ben jösz-hoz-zánk, mind-ö--rök-ké Ál-dott! REF. An-gyal-se-re-gek a menny-ben mind-nyá-jan ma-gasz-tal-nak Té-ged, s_ve-lük a_halandó em-ber, és min-den te-remt-mé-nyed di-csér. REF. A zsi-dók-nak né-pe pál-mák-kal jött az ú-ton e-léd, mi_is_könyörgé-sek-kel el-jöt-tünk í-me é-nek szó-val. REF. Mi-dőn_szenvedni men-tél, hó-dol-va di-csér-tek ők Té-ged, most,_midőn_mennyben tró-nolsz, mi zeng-jük e-lőt-ted a him-nuszt. REF. Tet-szett se-re-gük né-ked, tes-sék ma ne-ked szen-telt né-ped, mert_kegyelmes_ki-rály vagy, s_min-den jót tet-szé-sed-del fo-gadsz. REF.}
\end{gregosheet}

\commentary{A pap az oltárhoz megy, a szokott módon tiszteletadást végez, és esetleg megtömjénezi.
Majd minden más bevezető részt elhagyva, a Könyörgéssel kezdi a misét, és a szokott módon folytatja.}

\subsection{Mise}

\parttitle{Könyörgés}

\p Mindenható, örök Isten, Üdvözítőnk példát adott nekünk az alázatosságra, amikor szent
akaratodból emberi testet öltött és vállalta a kereszthalált. Segíts jóságosan, hogy
kínszenvedésének tanítását megértsük, és társai lehessünk feltámadásában. Aki veled él és uralkodik
a Szentlélekkel egységben, Isten mindörökkön-örökké.

\h Ámen.

\parttitle[Iz 50,4-7]{Olvasmány}

\lec Olvamány Izajás próféta könyvéből

Izajás így jövendölt a Megváltóról: Ezt mondja az Úr Szolgája: „Isten, az Úr tanította nyelvemet,
hogy az igével támasza lehessek a megfáradtaknak. Reggelenként ő teszi figyelmessé fülemet, hogy rá
hallgassak, mint a tanítványok. Isten, az Úr nyitotta meg fülemet. S én nem álltam ellen és nem
hátráltam meg. Hátamat odaadtam azoknak, akik vertek, arcomat meg azoknak, akik tépáztak. Nem
rejtettem el arcomat azok elől, akik gyaláztak és leköpdöstek. Isten, az Úr megsegít, ezért nem
vallok szégyent. Arcomat megkeményítem, mint a kőszikla, s tudom, hogy nem kell szégyenkeznem.

\lec Ez az Isten igéje.

\h Istennek legyen hála!

\parttitle[Tenuisti manum]{Graduále}

\begin{gregosheet}
  \melody{<ô-Ÿ---------------2--1---1r-4--:-£---------------------------2-----3r-5--,-£-----------------------2--1-----2r-4-.}
  \lyrics{   Te_tartottad_az én job-bo-mat, és_a_te_akaratod_szerint_ve-zetsz en-gem, és_a_te_dicsőségedbe_be-fo-gadsz en-gem.}
\end{gregosheet}

V. Mily jó Izráel Ist\underline{e}ne † az egyenes szívű\underline{e}khez, * az én lábaim mégis csaknem megtán/tor\underline{o}dtak. TE TARTOTTAD... \\
V. ! M\underline{e}rt megbotránkoztam a bűnösök m\underline{i}att, * látván a bűnösök bé/kess\underline{é}gét. TE TARTOTTAD...

\parttitle[Fil 2,6-11]{Szentlecke}

\lec Szentlecke Szent Pál apostolnak a filippiekhez írt leveléből

Testvéreim! Krisztus Jézus, mint Isten, az Istennel való egyenlőséget nem tartotta olyan dolognak,
amelyhez feltétlenül ragaszkodnia kell, hanem szolgai alakot öltött, kiüresítette önmagát, és
hasonló lett az emberekhez. Megalázta önmagát, és engedelmes lett a halálig, mégpedig a
kereszthalálig. Ezért Isten felmagasztalta őt, és olyan nevet adott neki, amely fölötte van minden
névnek, hogy Jézus nevére hajoljon meg minden térd a mennyben, a földön, és az alvilágban, és
minden nyelv hirdesse az Atyaisten dicsőségére, hogy Jézus Krisztus az Úr.

\lec Ez az Isten igéje.

\h Istennek legyen hála!

\parttitle[Deus, Deus meus]{Traktus}

\begin{gregosheet}
  \melody{<X-3--4--¥--------------------------:¥----------------7--rR3-:-6-----6----6----6---7--4--uU6-zZ5-4---3r-4--,,-¥--------------------------------6---6---7----rR3--,-4--6--6-/-6--7---4--uU6zZ5T4-3r-4-,,-8--8---8--8-8-----8--6----iI7-7uJ4R3-,-3--4--6--6---7--uJ4uU6zZ5T4-3r-4-,,-8--8---8---8---6---8-7---7uJ4R3-,-3---4---6--7---7---4--7--6---zZ5T4-3r-4--,,-¥----------------------------7--rR3-3--,-¥-------------------5---4---3---4---ed13r-tT4-,,-¥--------------------------------7--rR3-3--,-¥----------------:6--5-4---3---4--ed13r-tT4-,,-5----6---7----¦--------------------------------6---7--rR3-,-¥---------------------------6--6---5---4--3r--4--,,-5z--¦-------------7----6--7---rR3-,-6---6--¥---------:6--5---4--3r-4--,,-©----------------------6---iI7-u7J4R3-,-¥---------------------------------:¥-------------uJ4-uU6zZ5-4--3r--4--,,-©---------------------------6---8---7--7uJ4R3-,,-¥---------------¥---------------7---4-uU6zZ5T4-3r--4--,,-¥----------------------:6--6----7--rR3-3-,-¥----------------:7--4---7--6zZ5T4-3r--4-,,-6---6---¥------------------------6---6--6---7--rR3-3--,-3---4---¥------------------4---6zh4R3¨tT456u-rR3-,-4-6---6-/--6----6---6---6---5---B778ol7-uU6-;-6uU6-zZ5T4R3¨4tT4t¨uU6uU6h4t-/-8---8---7--8ik6Z5T4rR34z6u¨iI7j5z65z¨iI7U6-zZ5T4-tT4-.}
  \lyrics{   Is-te-nem,_Istenem,_tekints_reám, miért_hagytál_el en-gem?  Miért vagy mesz-sze ki-ál-tá--som sza-vá-tól?  Én_Istenem,_napestig_kiáltok,_és meg nem hall-gatsz? És éj-jel és nem fi-gyelsz   re-ám?  Pe-dig te a szent he-lyen la--kol,     Iz-ra-el-nek di-cső---------sé-ge!  Te-ben-ned bíz-tak a-tyá-ink,     bíz-tak és meg-sza-ba-dí-tot-tad   ő--ket.  Hozzád_kiáltottak,_és_megsza-ba-dul-tak, Tebenned_bíztak,_és meg nem szé-gye-nül---tek.   Én_pedig_féreg_vagyok_már_és_nem is em--ber, emberek_gyalázata és a nép meg-ve-tett--je.    Mind-nyá-jan, akik_látnak_engem,_megcsúfol-nak en-gem;      ajkukkal_megszólanak,_és_fe-jü-ket haj-to-gat-ják.  „Az Úrban_bízott, ment-se meg őt,   sza-ba-dítsa_meg, ha ked-ve-li őt!”  Csak_néznek_és_szemlél-nek en--gem,     elosztották_maguk_között_ruháimat, és_köntösömre sor-sot    ve-tet-tek. Szabadíts_meg_engem_az_orosz-lán szá-já-ból,      és_a_vadak_táma-dásá_tól_engem, meg-a-lá-------zot-tat.  Akik_félitek_az_U_rat, di-csér-jé-tek Őt, Jákobnak_ivadékai, di-cső-ít-sé-----tek Őt!  Hir-des-sétek_az_Urat_a_jö-ven-dő nem-ze-dék-nek, hir-des-sék_az_egek_az_Ő_i-gaz-sá------------gát   a nép-nek, mely egy-kor meg-szü-le------tik,  me---lyet                      maj-dan az Úr                              te----remt!}
\end{gregosheet}

\parttitle{Passió}

\commentary{Máté passió a Cantus Cantorumból.}

\parttitle{Homília}

\commentary{A szenvedéstörténet felolvasása után – ha alkalmas – mondható rövid homília.}

\parttitle{Felajánlási ének}

\begin{lilypond}
\layout {
  \context {
    \Score
    \omit BarNumber
  }
  \context {
    \Staff
    \omit TimeSignature
  }
}
\relative {
  \key d \major
  a'4 b8( a8) g4 b4 | a4( g4) fis2 | a4 a8( g8) fis4 e4 | d2. r4 | \break
  a'4 b8( a8) g4 b4 | a4( g4) fis2 | a4 a8( g8) fis4 e4 | d2. r4 | \break
  e4 e4 a2 | g4 fis4 e2 | e4 a4 g4 fis4 | e2. r4 | \break
  a4 b8( a8) g4 b4 | a4( g4) fis2 | a4 a8( g8) fis4 e4 | d2. r4 \fine
}
\addlyrics {
  \override LyricText.font-name = #"Cambria"
  \override LyricText.font-size = #0.1
  Meg -- vál -- tó ki -- rá -- lyunk e -- lé -- be me -- gyünk,
  mél -- tó tisz -- te -- let -- tel U -- runk -- hoz le -- gyünk!
  Ho -- zsan -- nát ki -- ált -- son né -- ki min -- den hív,
  di -- csé -- ret -- tel ál -- dást mond -- jon nyelv és szív!
}
\end{lilypond}

\begin{hymnstanzas}
  Így t\underline{i}sztelte J\underline{é}zust a zs\underline{i}dó nemzet,\\
  vir\underline{á}gvasárn\underline{a}pján nagy gy\underline{ü}lekezet.\\
  Elébe kimentek pálmaágakkal,\\
  dics\underline{é}rő én\underline{e}kkel és v\underline{i}rágokkal.
\stanzabreak
  Hogy \underline{a} szent vár\underline{o}shoz már k\underline{ö}zelgetett,\\
  és f\underline{é}nyes pomp\underline{á}val bév\underline{e}zettetett,\\
  ruháját a népség földre teríté,\\
  és \underline{ú}tját zöld \underline{á}ggal föl\underline{é}kesíté.
\end{hymnstanzas}
\centerstanza[0.5]{%
  És \underline{ö}rvendő sz\underline{í}vvel ki\underline{á}ltozának:\\
  Hozs\underline{á}nna, hozs\underline{á}nna Dáv\underline{i}d Fiának!\\
  Áldott az, ki az Úr szent nevében jön!\\
  Áldj\underline{á}k őt a h\underline{í}vek most \underline{é}s örökkön.
}

\parttitle{Felajánló könyörgés}

\p Urunk, Istenünk, egyszülött Fiad szenvedése engeszteljen ki téged irántunk, és amit
cselekedeteinkkel nem tudunk kiérdemelni, ennek a végtelen értékű áldozatnak érdeméből adja meg
nekünk a te irgalmad. Krisztus, a mi Urunk által.

\h Ámen.

\parttitle{Prefáció}

\p Az Úr legyen veletek. \h És a te lelkeddel.

\p Emeljük föl szívünket. \h Fölemeltük az Úrhoz.

\p Adjunk hálát Urunknak, Istenünknek. \h Méltó és igazságos.

\p Mert valóban méltó és igazságos, illő és üdvös, hogy mindig és mindenütt hálát adjunk néked, –
mi Urunk, szentséges Atyánk, mindenható örök Isten: Krisztus, a mi Urunk által.

Bár őmaga semmit sem vétett, az ítélet ártatlanul érte, – a bűnösökért mégis szenvedést vállalt. Az
ő halála bűneinket elvette, feltámadása megigazulást adott nekünk.

Az ő dicsőségét vallja ég és föld, vallják az angyalok és főangyalok, – hangos szóval hirdetik, és
vég nélkül zengik:

\parttitle{Sanctus}

\begin{gregosheet}
  \melody{<X-5-----tT4---:-5-----tT4-:-3-----4-----[--------------5--4--5-,-3--4---[-------------4-6----5----5--4-tT4-3--,-3--4----5--5-5--5---4---5--,-34t-5---/-5-5--5---4--6--5--4tT4-3--,-1--3r---4-!-4-4--5---rR3-3r-.}
  \lyrics{   Szent vagy, * szent vagy, szent vagy, mindenség_ura, Is-te-ne! Di-cső-séged_betölti a meny-nyet és a föl-det, Ho-zsan-na a ma-gas-ság-ban! Ál--dott, a-ki jön az Úr ne-vé---ben! Ho-zsan-na  a ma-gas-ság-ban!}
\end{gregosheet}

\parttitle{Agnus Dei}

\begin{gregosheet}
  \melody{<-5--5---5--rR3-4---:--4--[-------------3---4--eE2-1-,-3--4---5----4---4---,,-5--5---5--rR3-4---:-4--[-------------3---4--eE2-1-,-3--4---5----4---4---,,-5--5---5--rR3-4---:--4--[-------------3---4--eE2-1-,-3---4--5----rR3E2-2-.}
  \lyrics{  Is-ten Bá-rá--nya, * te elveszed_a_vi-lág bű-ne--it, ir-gal-mazz ne--künk!  Is-ten Bá-rá-nya, * te elveszed_a_vi-lág bű-ne--it, ir-gal-mazz ne--künk!  Is-ten Bá-rá-nya, *  te elveszed_a_vi-lág bű-ne-it,  adj ne-künk bé----két!}
\end{gregosheet}

\parttitle{Áldozási ének}

\begin{lilypond}
\layout {
  \context {
    \Score
    \omit BarNumber
  }
  \context {
    \Staff
    \omit TimeSignature
  }
}
\relative {
  \key es \major
  \time 3/4
  es'4 es4 es4 | es2 f4 | g2 g4 | f2. | \break
  f4 as4 g4 | es2 g4 | f2 f4 | es2. | \break
  g4 g4 g4 | g2 f4 | g2 as4 | bes2. | \break
  f4 as4 g4 | es2 g4 | f2 f4 | es2. \fine
}
\addlyrics {
  \override LyricText.font-name = #"Cambria"
  \override LyricText.font-size = #0.1
  A meg -- vál -- tó szent kín -- ha -- lál,
  mely -- ben a föld üd -- vöt ta -- lál,
  ad -- jon, U -- runk, most bé -- kü -- lést,
  szí -- vünk vá -- gyá -- nak eny -- hü -- lést!
}
\end{lilypond}

\begin{hymnstanzas}
  Tápláljanak szent titkaid,\\
  ihlessenek vércseppjeid,\\
  lelkünknek mélye teljen el\\
  szenvedésed emlékivel.
\stanzabreak
  Éljen szívünkben szüntelen\\
  e gyalázat, e gyötrelem,\\
  keresztfa, lándzsa, vas-szögek,\\
  tövisszúrások, sebhelyek!
\end{hymnstanzas}
\centerstanza{%
  Ó, megfeszített Jézusunk,\\
  egész szívvel fohászkodunk.\\
  Kit eladott a gyűlölet:\\
  Tenéked áldás, tisztelet.
}

\parttitle[Pater, si non potest]{Kommunió}

\begin{gregosheet}[tone=8. dó tónus]
  \melody{<XB-zZ5-4-----:-3--4--4---6---4---3-2rR3--3--:-5----7--6--5uU6-5---,-6--6uU6-zZ5z-:-4-5--6-4--5--3-,,-3-4-¥-†-¥-5-6-7-ˆ-6-,-¥-4-6-7-ˆ-6.}
  \lyrics{    A---tyám, * ha el nem múl-hat e po----hár, hogy ki ne i----gyam, le-gyen meg    a te a-ka-ra-tod!}
\end{gregosheet}

\begin{psalmsheet}[tone=8, continuous]
  \psalmtext{
    1. Istenem, Istenem, miért hagytál el engem * miért vagy messze kiáltásom szavától?\\
    2. Én Istenem, napestig kiáltok, és meg nem hallgatsz * és éjjel, és nem figyelsz reám?! ANT. ---

    3. Mert körülvettek engem az ebek * a gonoszok gyülekezete reám szállott.\\
    4. Átlyuggatták kezeimet és lábaimat * és megszámlálták minden csontomat. ANT. ---

    5. Szabadíts meg engem az oroszlán szájából * és a vadak támadásától engem, megalázottat!\\
    6. És én hirdetem a te nevedet atyámfiainak * az egyháznak közepette dicsérlek téged. ANT. ---

    7. Akik félitek az Urat, dicsérjétek őt * Jákobnak ivadékai, dicsőítsétek őt!\\
    8. Emlékeznek és az Úrhoz térnek a föld minden határai * és hódolnak előtte a népek családjai. ANT. ---

    9. És az én lelkem őneki él * és az én magzatom néki szolgál.\\
    10. És hirdettetik az Úr a jövendő nemzedéknek † hirdetik az egek az ő igazságát a népnek, mely egykor megszületik * melyet majdan az Úr teremt. ANT.
}
\end{psalmsheet}

\parttitle{Áldozás utáni könyörgés}

\p Mennyei eledellel tápláltál minket, Istenünk. Könyörögve kérünk, hogy akik szent Fiad halálából
már biztosan remélhetjük azt, amiben hiszünk, feltámadásának erejéből eljuthassunk oda, ahová
törekszünk. Krisztus, a mi Urunk által.

\h Ámen.

\parttitle{Kivonulási ének}

\begin{lilypond}
\layout {
  \context {
    \Score
    \omit BarNumber
  }
  \context {
    \Staff
    \omit TimeSignature
  }
}
\relative {
  \key bes \major
  \partial 4 d'4 | g4 f4 es4 d4 | c2 d4 \breathe a'4 | bes4 bes4 a8( g8) a4 | g2. \break
  \partial 4 d4 | g4 f4 es4 d4 | c2 d4 \breathe a'4 | bes4 bes4 a8( g8) a4 | g2. \break
  \partial 4 bes4 | a8( g8) f4 g4 a4 | bes2 bes4 \breathe f4 | g4 f4 es4 es4 | d2. \break
  \partial 4 bes'4 | a8( bes8) c4 bes4 a4 | g2 a4 \breathe d,4 | es4 d4 c4 f4 | d2. \fine
}
\addlyrics {
  \override LyricText.font-name = #"Cambria"
  \override LyricText.font-size = #0.1
  Szent ar -- cod fé -- nyes -- sé -- ge, jajj, hol van, Jé -- zu -- som?
  És hol van sze -- med fé -- nye? El -- rej -- ti fáj -- da -- lom.
  Nincs ben -- ned sem -- mi ép -- ség, ott hal -- do -- kolsz a fán,
  és tes -- ted csu -- pa kék -- ség az üt -- le -- gek nyo -- mán.
}
\end{lilypond}

Mint áldozati bárány, ha oltárr\underline{a} kerül,\\
te véred ontod némán, az Írás t\underline{e}ljesül.\\
Ó sz\underline{á}nj meg, Uram, nézd el sok balgaságomat,\\
megv\underline{a}llom szájjal, szívvel, hogy Isten Fia vagy!

\section{Nagycsütörtök}

\subsection{Esti mise az utolsó vacsora emlékére}

\parttitle[Nos autem gloriari]{Introitus}

\begin{gregosheet}[tone=1. lá tónus]
  \melody{<X-1--1--qQ0-3--4---3t--5----:-6-5-----4--4---4-----4---4--4-----3t-4--,-4---4--4---2-4--tT4-3---:-3-3--2----3--4---3--2--1--1--,,-3-4-[-4-3-:-[-4-3-†-rR3-1-.}
  \lyrics{   Mi pe-dig di-cse-ked-jünk * U-runk, Jé-zus Krisz-tus ke-reszt-jé-ben; ben-ne van a mi üd--vünk, é-le-tünk és fel-tá-ma-dá-sunk!}
\end{gregosheet}

\begin{psalmsheet}[tone=1, continuous]
  \psalmtext{
    1. Könyörüljön rajtunk az Isten, és áldjon meg minket, * világosítsa ránk az Ő orcáját, és könyörüljön rajtunk! ANT. \\
    2. Hogy megismerjük a földön a Te utadat, * minden nemzet meglássa a Te üdvözítő művedet! ANT. \\
    3. Hálát adjanak neked a népek, ó, Isten, * hálát adjanak neked minden nemzetek! ANT.
}
\end{psalmsheet}

\commentary{Nincs "Dicsőség az Atyának..."}

\parttitle[Cunctipotens]{Kyrie}

\commentary{A Kyriet váltakozó karokban énekeljük, minden sort háromszor. Az utolsó Kyrie eleison
végét közösen énekeljük. Vagy énekesek előénekelnek minden sort és utána közösen megismételjük
kétszer; az utolsó Kyrie eleison-t az előénekesek kezdik és közösen fejezzük be.}

\begin{gregosheet}
  \melody{<-5--tT4-5¨uU6Z5T4t-:-tT4R3d1w-:-2rt²1¨2rR3E2-3--1-1-,,-5-----5¨4tT4R3¨44tg3¨4tG2W1-:-2rR3E2-3--1-1-,,-1t-tT4-5uU6h4¨zZ5¨7uj5-:-5zZ5¨rR3E2-4--5-5-,,-1t-tT4-5uU6h4¨zZ5¨7uj5-:-1t¨tT4¨5uU6h4¨zZ5¨7uj5-:-5zZ5¨rR3E2-4--5-5-.}
  \lyrics{  Ky-ri--e,           @          e------------le-i-son! Chris-te                      e------le-i-son! Ky-ri--e                 e----------le-i-son! Ky-ri--e                 @                      * e----------le-i-son!}
\end{gregosheet}

\parttitle[Missa mundi]{Gloria}

\vskip\blockvskip

\p

\begin{gregosheet}
  \melody{<-2--4---5---5-5--5---4---5---4--2---2-,}
  \lyrics{  Di-cső-ség a ma-gas-ság-ban Is-ten-nek!}
\end{gregosheet}

\commentary{Ezután meghúzzák a harangokat a templomban található minden csengővel folyamatosan
csöngetnek, az orgona prelúdiumot játszik.}

\vskip\blockvskip

\h És a földön békesség a jóakaratú embereknek! Dicsőítünk Téged, áldunk Téged, imádunk Téged,
magasztalunk Téged. Hálát adunk Neked nagy dicsőségedért. U\-runk és Istenünk mennyei Király,
mindenható Atyaisten. Urunk Jézus Krisztus, egy\-szülött Fiú. Urunk és Istenünk, Isten Báránya, az
Atyának Fia, Te elveszed a világ bűneit, irgalmazz nekünk! Te elveszed a világ bűneit, hallgasd meg
könyörgésünket! Te az Atya jobbján ülsz, irgalmazz nekünk! Mert egyedül Te vagy a Szent, te vagy az
Úr, te vagy az egyetlen Fölség, Jézus Krisztus, a Szentlélekkel együtt az Atyaisten dicsőségében.
Ámen.

\commentary{Ezutána a harangok és csengők elhallgatnak egészen a húsvéti vigíliáig.}

\parttitle{Könyörgés}

\p Istenünk, annak a szent vacsorának megünneplésére jöttünk össze, amelyen egyszülött Fiad, új és
örökre szóló áldozatát és szeretetlakomáját halála előtt az Egyházra bízta. Kérünk, segíts, hogy
ebből a nagy misztériumból a szeretet és az élet teljességét meríthessük. A mi Urunk, Jézus
Krisztus, a te Fiad által, aki veled él és uralkodik a Szentlélekkel egységben, Isten
mindörökkön-örökké.

\h Ámen.

\parttitle[Kiv 12,1-8.11-14]{Olvasmány}

\lec Mózes második könyvéből

Az Úr így szólt Mózeshez és Áronhoz Egyiptomban: „Ez a hónap legyen számotokra a kezdő hónap: ez
legyen az év első hónapja. Hirdessétek ki Izrael egész közösségének: A hónap tizedik napján
mindenki szerezzen egy bárányt családon\-ként, egy bárányt házanként. De ha a család kicsi egy
bárányhoz, akkor a személyek számának megfelelően a szomszédos családdal együtt vegyen egyet.
Aszerint számoljátok a meghívottakat, hogy ki-ki mennyit eszik. Az állat legyen hibátlan, hím és
egyéves. Vehettek bárányt vagy kecskét. Tartsátok a hónap tizennegyedik napjáig. Akkor Izrael
közösségének egész gyülekezete estefelé vágja le. Vegyenek a véréből és kenjenek belőle annak a
háznak a két ajtófélfájára és szemöldökfájára, amelyben elköltik. A húsát – tűzön megsütve – még
akkor éjszaka egyék meg. Kovásztalan kenyérrel és keserű salátával fogyasszák. Így fogyasszátok:
Legyen a derekatok felövezve, saru a lábatokon, bot a kezetekben. Sietve egyétek, mert ez az Úr
átvonulása. Én végigvonulok azon az éjszakán Egyiptomon és megölök minden elsőszülöttet
Egyiptomban: embert és állatot, s ítéletet tartok Egyiptom minden istene fölött, én, az Úr. A vért
használjátok annak a háznak a megjelölésére, amelyben laktok. Ha látom a vért, kihagylak
benneteket. Titeket nem ér a megsemmisítő csapás, amellyel Egyiptomot megverem. Ez a nap legyen
számotokra emléknap, és üljétek meg úgy, mint az Úr ünnepét. Nemzedékről nemzedékre tegyétek meg
ünnepnapnak mindörökre.”

\lec Ez az Isten igéje.

\h Istennek legyen hála!

\parttitle[Kiv 12,1-8.11-14]{Graduále}

\begin{gregosheet}
  \melody{<X-33r---3-:-¢----------------------4----3--4t-tT4-,-4-3--4-----5--rR3-3-,,-[--------------------------5z-5--:-[--------------4t-4-:-£----------------5---rR3-4t--tT4-.}
  \lyrics{   Krisz-tus engedelmes_volt_értünk mind-ha-lá-lig,  a ke-reszt-ha-lá--lig. Azért_felmagasztalta_őt_az Is-ten, és_nevet_adott né-ki, mely_fölötte_van min-den név-nek.}
\end{gregosheet}

\commentary{Vagy díszesebb dallamon:}

\begin{gregosheet}
  \melody{<X-33r---3-:-¢---------------4tg3E2¨11rf2¨33r---3-:--3t34tu7¨iI7i™uU6u55¨44-4-5--7uj5g3¨zh4tT4t¨3rR3-3-,--3-3--5z----445z¨uU6h44t¨zZ5T4-3---3---3---eD0¨'3r¨tT4R3E24t3tT4t¨rR3-,,--3-3----3--tT4-4u--7-(-¦--------------iI7uJ4-:-7iI7iI7uJ4¨!7iöõÕö8ouj5¨™uU6¨8ol7oO8I7uU6Z5-,--[--------5-7i---iI7¨9pé8I7-uj5-:-¢----------------------3tT4R3¨4t¨™7uU6Z5-uj5¨XzZ5zh4R3¨5zh4R3¨tT477u¨5t¨3tT4t¨rR3.}
  \lyrics{   Krisz-tus engedelmes_volt ér-----------------tünk mind         @         a ha-lá------------------lig, a ke-reszt-nek                ha-lá---lá--ig.  @                         A-zért az Is--ten is, felmagasztalta Őt,      @           @            @                     és_nevet a-dott né---------ki,   mely_minden_név_fölött va---------@------gyon.}
\end{gregosheet}

\parttitle[1Kor 11,23-26]{Szentlecke}

\lec Szentlecke Szent Pál apostolnak a korintusiakhoz írt első leveléből

Én az Úrtól kaptam, amit közöltem is veletek: Urunk Jézus azon az éjszakán, amikor elárulták, fogta
a kenyeret, hálát adott, megtörte, és így szólt: „Vegyétek és egyétek, ez az én testem, amely
értetek adatik! Ezt cselekedjétek az én emlékezetemre!” A vacsora után ugyanígy fogta a kelyhet is,
és így szólt: „Ez a kehely az újszövetség kelyhe az én véremben. Valahányszor isztok belőle,
tegyétek ezt az én emlékezetemre!” Amikor ugyanis e kenyeret eszitek és e kehelyből isztok, az Úr
halálát hirdetitek, amíg el nem jön.

\lec Ez az Isten igéje.

\h Istennek legyen hála!

\parttitle[In hoc cognoscent]{Traktus}

\commentary{Egy szólista énekli a következő traktust:}

\begin{gregosheet}
  \melody{<-4--tT4-4u-¦-------------------------------------tT4u-uj5u¨rR3-,-3t---¦--------------------------7-7uU6Z5¨zZ5T4¨tT4-,-4--tT4-4u--¦--------------7i7j5-5-:-¦--------------------7---tT4u-uj5u¨rR3,-3t-¦-----------------7-7---7---7uU6Z5¨zZ5T4¨tT4-,-4----tT4-4u--7--7---7-7iI7j5-5--:-¦-----------7----7--7-7---tT4u-uj5u¨rR3-,-3t-¦-------------------------7uU6Z5¨zZ5T4¨rR3-,-4---tT4-4u--7--7---7iI7j5-5-:--¦------------------------tT4u-uj5u¨rR3-,-3---5--¦------------------------------------77¨iI7U6Z5T45uU6Z5¨4t¨33¨55z¨uU6Z5z¨55¨44-.}
  \lyrics{  Ar-ról is-merje_meg_mindenki,_hogy_tanítványaim vagy-tok        hogy szeretettel_vagytok_egymás i-ránt.              Ma-rad-jon meg_bennetek_e há----rom a_hit,_a_remény_és_a sze-re---tet       de legnagyobb_köztük a sze-re-tet.                Most még meg-ma-rad e há-----rom: a_hit,_a_re-mény és a sze-re---tet        de legnagyobb_köztük_a_szere-tet.               Sze-res-sük te-hát egy----mást mert_a_szeretet_Istentől va-ló           s_a-ki szereti_testvérét,_az_Istentől_szüle-tett.}
\end{gregosheet}

\end{document}
